% status: FULL
\chapter{Model File Formats: GGUF and safetensors}\label{app:formats}
\begin{ChapterIntent}
Loading a checkpoint is an interpretation. A loader must agree with the
model on names, shapes, encodings and ordering. This chapter separates a
valid container from a compatible model and supplies an address calculation
to audit by hand. We will also construct a 256-byte GGUF file, decode every
weight and execute two projections through the real ggml CPU backend.
The file is a numerical teaching fixture, not a miniature trained LLM.
\end{ChapterIntent}

\section{Three contracts}
The container locates bytes. The tensor contract interprets them as shaped
elements. The architecture assigns tensors to operations, including the
tokenizer and position convention. Passing the first contract does not
establish the other two. A correctly bounded array can still be the wrong
projection matrix.

Inspect metadata, validate sizes and required names, check encodings,
construct the model, and run a known-input comparison before timing.
Record the checkpoint's content hash and configuration with the engine
revision. A mutable model alias does not identify immutable bytes.

\section{safetensors offsets by hand}
The format starts with an eight-byte little-endian header length, followed
by a UTF-8 JSON header and a data buffer. Tensor offsets are relative to
the buffer, with an exclusive end. Entries declare dtype and shape;
ordering is row-major. The format avoids pickle's arbitrary object
execution, but numerical values can still be nonfinite
\cite{safetensors2026format}.

For an illustrative 120-byte header, data begins at byte 128.
An F32 tensor of shape $[2,3]$ and offsets $[0,24)$ occupies absolute
bytes $[128,152)$. Six four-byte elements account for that interval.
Element $(1,2)$ has linear index $1\cdot3+2=5$ and begins at byte
$128+5\cdot4=148$ (derived).

\EngineFigure{foundation-file-offsets}{Illustrative byte intervals with
exclusive ends. The data origin must be added exactly once. The declared
header length includes its padding, not the eight-byte length field.
The drawing is schematic, not proportional to byte counts.}{fig:foundation-file-offsets}

\section{GGUF and quantized blocks}
GGUF carries a versioned header, key--value metadata, tensor descriptors,
alignment padding and tensor data. Descriptors specify dimensions, type
and an offset relative to the tensor-data region. Alignment is part of
the format, not permission to round an invalid interval back into bounds
\cite{ggml2026gguf}.

A quantized type describes packed blocks and their metadata. Its size is
not generally element count times an integer bytes-per-element value.
Use the declared block geometry and byte size, check required divisibility,
and account for mixed tensor types. A four-bit checkpoint label is not
an exact file-size formula. Chapter~\ref{app:numerics} shows why scales
change the effective bit count.

\section{Read the first twenty-four bytes}
Begin with a smaller question than loading a language model: what can a
reader establish before it allocates a single tensor? Our fixture starts
with the following hexadecimal bytes. A pair of hexadecimal digits
represents one byte; two pairs are not automatically one integer.
\begin{verbatim}
47 47 55 46  03 00 00 00
02 00 00 00 00 00 00 00
01 00 00 00 00 00 00 00
\end{verbatim}
The first four bytes spell \texttt{GGUF}. They identify the container,
not its model architecture. The next four encode the integer three in
little-endian order. For a four-byte unsigned field with bytes
$b_0,b_1,b_2,b_3$, the value is
\begin{equation}
u=b_0+2^8b_1+2^{16}b_2+2^{24}b_3.
\end{equation}
Each $b_i$ is an integer from zero to 255. The byte nearest the beginning
of this field carries the smallest positional weight. Consequently
\texttt{03 00 00 00} means three, while \texttt{00 00 00 03} would
mean $3\cdot2^{24}$ under this interpretation. Endianness changes how
bytes make a scalar; it does not reverse a tensor's logical axes.

The next two fields are eight-byte unsigned counts: two tensors and
one metadata entry. They end at absolute byte 24. Offsets in this chapter
are zero-based and intervals exclude their upper endpoint. Thus the
version occupies $[4,8)$, not bytes four through eight inclusive.
That convention removes a common one-byte ambiguity from bounds checks.

\begin{center}
\begin{tabular}{@{}rrll@{}}
\toprule
Start & Bytes & Field & Fixture value\\
\midrule
0 & 4 & Magic bytes & \texttt{47 47 55 46}\\
4 & 4 & Version & 3\\
8 & 8 & Tensor count & 2\\
16 & 8 & Metadata count & 1\\
\bottomrule
\end{tabular}
\end{center}
The version-three layout is defined by the GGUF specification
\cite{ggml2026ggufpinned}. The lab deliberately accepts only its
little-endian subset. It rejects other version encodings; it does not
claim that every rejected file is invalid GGUF. Supporting another
byte order means interpreting every affected field and payload
consistently, not swapping the version and hoping the rest works.

Do not map the file directly onto an ordinary compiler-generated C
structure. A compiler may insert padding between members, and the
variable-length regions do not fit fixed-size members in the first
place. A serialized layout is a sequence of explicit fields. Its parser
advances by the width actually specified for each field.

\section{Metadata is typed data, not decoration}
At byte 24 we store a string naming the key
\texttt{general.alignment}. A GGUF string starts with an eight-byte
byte count, then that many UTF-8 bytes. This key has 17 ASCII bytes,
so its complete string representation needs $8+17=25$ bytes and ends
at byte 49. There is no terminating zero byte to skip.

The next four bytes identify the metadata value type. Here the type
code is four, meaning \texttt{UINT32}. The following four bytes
encode the value 32. The entry therefore occupies
\[
8+17+4+4=33\text{ bytes},
\qquad 24+33=57.
\]
We now know where the tensor directory starts without inspecting any
tensor data. If a parser treated the string length as a four-byte field,
every later read would be shifted. Some shifted values might still look
plausible; plausibility is not structural validation.

The type tag is not optional information. A string containing the
characters \texttt{32} and an unsigned integer with value 32 have
different serialized lengths and different interpretations. Neither
should silently substitute for the other. Likewise, the metadata
type code and the tensor type code belong to different enumerations.
Code four in a metadata field is not permission to decode a tensor
with type four. Our supported tensor types are F32, code zero, and
Q4\_0, code two. The tests deliberately swap these namespaces.

String lengths count bytes, not displayed characters. The optional
metadata test stores \texttt{café}; its UTF-8 encoding has five bytes.
Another test stores a three-byte scalar in a token array. The reader
must consume the declared byte span before decoding the text. Advancing
by a programming language's character count can leave the cursor inside
a multibyte scalar, corrupting the next field.

An array has another distinction: its length counts elements, while
the size of each element depends on its declared type. An array of
strings therefore needs a string-length field for every element.
Multiplying the array length by eight would account only for those
length fields, not for the strings themselves. The lab supports bounded
arrays of strings; it rejects nested arrays and other array element
types rather than guessing their widths.

This is where a loader can fail before doing any model arithmetic.
A damaged count can otherwise ask it to allocate billions
of entries or spend minutes walking an impossible directory. Bound
the count before entering the loop. Bound each length against the
remaining input before slicing. A small file should not create an
unbounded amount of parser work.

The reader uses one cursor for every field. These are the actual lab
methods, with line numbers matching the companion source:
\par\medskip\noindent\begin{minipage}{\linewidth}
\CodeLines{../code/foundations/gguf_bytes.py}{30}{47}
\end{minipage}\par\medskip
The subtraction in \texttt{take} asks how many bytes remain before
advancing the cursor. The format character \texttt{Q} requests an
eight-byte unsigned integer; \texttt{I} requests four bytes. The
\texttt{<} prefix fixes little-endian interpretation independently of
the machine running Python. A string consumes its bytes once and
then checks UTF-8. Neither a failed decode nor an excessive length
returns a partly interpreted field.

\section{A tensor descriptor locates a representation}
Our first tensor is named \texttt{dense.weight}. Its descriptor
contains the name, rank two, dimensions $[4,2]$, F32 type and relative
offset zero. Its name has twelve bytes. The descriptor needs
\[
(8+12)+4+(2\cdot8)+4+8=52\text{ bytes}.
\]
It starts at 57 and ends at 109. The second descriptor names
\texttt{packed.weight}; its thirteen-byte name makes its descriptor
one byte longer. It occupies $[109,162)$ and declares dimensions
$[32,1]$, Q4\_0 type and relative offset 32. The directory ends at
162. No weight has been read yet.

\EngineFigure{foundation-gguf-bytes}{The complete 256-byte teaching
container. Intervals are exact, but panel widths are schematic. The
descriptor offsets zero and 32 are relative to the data origin at 192,
not to the file's beginning. Unused bytes are explicit alignment
padding, not extra weights.}{fig:foundation-gguf-bytes}

The directory describes a logical tensor and a physical encoding at
the same time. The logical dimensions describe how many decoded
values exist along each axis. The type determines how those values
occupy bytes. With F32, eight values occupy 32 bytes. With Q4\_0,
32 values occupy one 18-byte block. Elements and stored bytes
are therefore not interchangeable units.

For a supported block format, let $n_0,\ldots,n_{r-1}$ be the
descriptor dimensions, $q$ the number of logical elements per block,
and $s$ the stored bytes per block. In the contiguous layout used here,
\begin{equation}
N_{\mathrm{bytes}}=
\left(\frac{n_0}{q}\right)
\left(\prod_{i=1}^{r-1}n_i\right)s,
\qquad n_0\bmod q=0.
\label{eq:gguf-block-bytes}
\end{equation}
The fastest dimension must contain whole blocks. Checking only that
the total element count is divisible by $q$ is insufficient: two rows
of sixteen elements have 32 elements in total, but neither row contains
a complete Q4\_0 block under this contract. A loader must not pack
across a row boundary unless the representation explicitly permits it.

Our parser keeps the offset field's own location in the returned
descriptor. It is byte 101 for the first tensor and 154 for the second.
Those are useful places to introduce a controlled fault in a test.
They are not addresses of the tensor data. The fact that one file
contains both an offset field and the location it describes is why
annotated byte diagrams are worth drawing.

\section{Alignment gives the data region its origin}
The directory ends at 162, but the declared alignment is 32 bytes.
The data region starts at the next aligned position:
\begin{equation}
D=A\left\lceil\frac{E}{A}\right\rceil
=32\left\lceil\frac{162}{32}\right\rceil=192\text{ bytes}.
\end{equation}
Here $E$ is the directory-end offset, $A$ the alignment in bytes and
$D$ the data origin. The 30 bytes in $[162,192)$ are padding.
A boundary already divisible by $A$ needs no additional
padding. Unconditionally adding an alignment unit would move an
already aligned region too far.

If a tensor's descriptor offset is $o$, its absolute payload begins
at $D+o$. For the two tensors this is 192 and 224. The F32 payload
ends at 224; the Q4\_0 payload ends at 242. Our canonical fixture
then pads the last allocation to 256. Its total size is therefore
24 fixed header bytes, 33 metadata bytes, 105 descriptor bytes,
30 directory-padding bytes, 50 weight bytes and 14 final-padding bytes.
Adding these independently gives 256.

There are three different sizes to report. The file is 256 bytes.
Its useful tensor payload is 50 bytes. Its two aligned tensor
allocations occupy 64 bytes. None of these is a contradiction.
They answer different questions: how much to transfer as a file,
how much represents weights, and how much this serialized allocation
layout reserves. Device repacking can introduce a fourth size.

For production code, alignment deserves more care than a memorized
bit trick. Rounding with \texttt{(x+A-1) \& \textasciitilde(A-1)}
requires a power-of-two alignment and an addition that does not
overflow. A general modulo-based calculation avoids the power-of-two
assumption, but its result still needs bounds checking. Python's
arbitrary-sized integers avoid machine wraparound; they do not remove
the need to bound resource use.

There is a concrete specification-versus-implementation boundary here.
The pinned GGUF document describes alignment as a multiple of eight.
The pinned \texttt{gguf.cpp} reader checks for a power of two
\cite{ggml2026ggufpinned,ggml2026loaderpinned}. For example, 24 satisfies
the first rule but not the second. Our lab chooses the intersection
and limits it to powers of two from eight through 4096. Rejecting 24
is a stated compatibility policy, not a claim that the two upstream
descriptions say the same thing. Record the implementation revision
when a format detail affects portability.

\section{Validate before allocating}
Metadata can claim enormous dimensions or offsets near integer limits.
Use checked arithmetic for products, additions and conversions. Validate
intervals against actual buffer length before constructing views. Reject
duplicate names and unsupported types; apply the format's overlap and
coverage rules. One parser silently overriding an entry another parser
uses is not a safe interpretation contract.

Shape checks should name axes. Equal element counts do not make
$[d,m]$ and $[m,d]$ interchangeable. Also check tied versus untied output
weights, normalization parameters, head counts and position conventions.
A well-formed file can still generate entirely wrong answers.

Configuration and preprocessing are part of identity. Two containers
holding equal tensor bytes do not prove equal inference when tokenization,
special tokens, chat templates or positional scaling differs. Keep these
dependencies in the reproduction record, not only in a filename.

\section{Mapped is not resident}
Memory mapping associates contents with virtual addresses. It does not
prove pages are resident, on an accelerator, or free from first-touch
faults. A warm run may benefit from the OS cache after a model process
exits. A mapping can retain its address while physical pages are reclaimed.

Distinguish file bytes, mapped address space, resident host bytes and device
allocations. Include repacking and conversion buffers in peak memory.
Do not mutate a shared checkpoint to make a benchmark easier.
Chapter~\ref{ch:offloading} follows these transfers; Chapter
\ref{app:reproducing} explains how to state warm and cold conditions
without claiming a cache flush that was never performed.

\section{Build a file whose every byte has a job}
The worked offset example is not only a drawing. The format lesson constructs
the file in memory and reads it back:
\begin{verbatim}
python3 code/foundations/lesson.py formats
\end{verbatim}
Its six F32 elements are the numbers one through six. The program creates
a JSON descriptor for a tensor named \texttt{toy}, pads the header with
spaces to 120 bytes, prefixes its eight-byte length and appends the six
little-endian values. It does not save or modify any model file.

Before running it, write down three predictions: the data origin is 128,
the absolute interval ends at 152, and element $(1,2)$ is 6. The output
confirms all three. The third prediction tests interpretation, not only
arithmetic: a wrong origin or wrong element width can produce an in-bounds
read of the wrong value.

Open \texttt{tensor\_file} and \texttt{inspect\_tensor} in
\texttt{code/foundations/lesson.py}. The inspector deliberately supports
only this one-tensor fixture. It rejects another shape, another dtype,
truncated data and unexplained trailing bytes. A small teaching parser
with explicit limitations is easier to audit than a supposedly general
loader whose untested branches accept arbitrary inputs.

Notice the order of checks. It checks that the length field exists before
reading it. It checks the declared header length against both a policy
limit and the actual file length before parsing JSON. Only then does it
read shape and interval metadata. Finally it checks the payload size before
interpreting floats. Each step establishes the precondition of the next.

\section{A shape is more than an element count}
Suppose a projection expects four input features and three output
features. The mathematical row-vector matrix has shape $[4,3]$.
A library that stores one row per output may store the same map as
$[3,4]$. Both contain twelve values. Reinterpreting one layout as the
other does not transpose the values; it changes which values multiply
which input features.

Test the loader with an input whose coordinates are deliberately
different, such as $[1,2,4,8]$. A row of identical inputs can hide a
permutation. Compare the loaded projection to a hand-defined map before
trying an entire model. The smaller boundary tells you whether the
error lies in tensor interpretation or later computation.

The same principle applies to fused projections. A checkpoint may store
Q, K and V separately, or concatenate them along a declared axis.
The engine must know where each region begins and how heads are ordered.
Dividing the total element count by three is wrong for grouped-query
attention when query width differs from key and value width.
The model configuration and the tensor names jointly determine the split.

An embedding table and an output matrix may be tied. In that case one
underlying parameter array plays two roles. Duplicating it can waste
memory; independently modifying one copy changes the model. Conversely,
assuming a tie where the checkpoint stores distinct arrays substitutes
different parameters. The loader is where that distinction becomes
an executable ownership decision.

\section{Read a quantized tensor as blocks, not small floats}
A packed integer code does not have its own universal real value.
Decoding usually needs a scale and sometimes a zero point or other
metadata associated with its block. The shape describes the logical
array after decoding; the byte interval describes the representation
before decoding. These are different coordinate systems.

For a hypothetical block of 32 four-bit codes and one 16-bit scale,
the representation occupies 18 bytes. Two such blocks encode 64
logical values in 36 bytes. Multiplying 64 by half a byte predicts
only the 32-byte code payload and omits the scales. This is a teaching
format, not a definition of any named GGUF quantization type.
For a real type, use its exact block structure.

Conversion also has a peak-memory cost. If a loader retains packed
weights while creating a second, unpacked array, both can be live at
once. A model that fits after conversion can fail during conversion.
Write the lifetime of the source, temporary workspace and destination
before computing the peak. Memory mapping does not automatically
eliminate that overlap.

\subsection{Decode an actual Q4\_0 block}
We can now replace the hypothetical block with a named representation.
At the pinned ggml revision, \texttt{block\_q4\_0} contains a binary16
scale and sixteen packed bytes. The decoder assigns low nibbles to
the first sixteen positions and high nibbles to the next sixteen
\cite{ggml2026q4pinned}. Our fixture's block is
\begin{verbatim}
00 38  f0 e1 d2 c3 b4 a5 96 87
78 69 5a 4b 3c 2d 1e 0f
\end{verbatim}
The first two bytes encode a binary16 scale of one half. As a
little-endian 16-bit word they form hexadecimal \texttt{3800}.
Its sign bit is zero, exponent field is fourteen, and fraction is
zero. With exponent bias fifteen, the normal-value formula gives
$2^{14-15}(1+0)=0.5$. The bytes are not two independent scales,
nor an F32 value truncated to half its bytes.

\EngineFigure{foundation-q4-nibbles}{One Q4\_0 block from the fixture.
The first packed byte is hexadecimal F0. Its low nibble describes
logical position zero; its high nibble describes position sixteen.
The separation into two halves is part of the representation.
This is decoding, not a recipe for choosing a quantization scale.}{fig:foundation-q4-nibbles}

Let $u_j$ be packed byte $j$ after the scale, for $0\leq j<16$.
Let $d$ be the decoded scale. The reconstruction is
\begin{align}
w_j &= d\bigl((u_j \mathbin{\&} 15)-8\bigr),\\
w_{j+16} &= d\bigl((u_j \gg 4)-8\bigr).
\end{align}
The bitwise AND keeps the low four bits; the right shift moves
the high four bits into the low positions. Each nibble is first
interpreted as an unsigned code from zero through fifteen.
Subtracting eight centers that code. These are offset codes,
not four-bit two's-complement integers.

For \texttt{F0}, the low nibble is zero and the high nibble is
fifteen. Therefore $w_0=0.5(0-8)=-4$ and
$w_{16}=0.5(15-8)=3.5$. The next byte, \texttt{E1}, gives
$w_1=-3.5$ and $w_{17}=3$. A decoder that alternates low and
high values would return $[-4,3.5,-3.5,3,\ldots]$ instead of
the correct first-half order. Both outputs contain the same
multiset of values; a check of their sum would miss the error.

This is why the reference check uses basis vectors. Multiplying by
the vector with a one at position $i$ and zeros elsewhere selects
weight $w_i$. Thirty-two such probes establish the coordinate order
in this one block. A single all-ones input checks only its sum,
which is negative eight. It cannot prove the positions are right.
Choose test inputs that distinguish the failure you are trying to find.

The complete block decoder is short enough to audit against the equation:
\par\medskip\noindent\begin{minipage}{\linewidth}
\CodeLines{../code/foundations/gguf_bytes.py}{155}{163}
\end{minipage}\par\medskip
The format character \texttt{e} decodes binary16. The two list
comprehensions preserve their own coordinate ranges, and
\texttt{low + high} concatenates those ranges. Replacing that
concatenation with alternating pairs would be an implementation error
even though the block's size and sum stayed correct.

Eighteen bytes for 32 values is $18\cdot8/32=4.5$ bits per value.
That is the block's effective storage rate, not four, and not the
whole file's compression ratio. Header, directory, padding and other
tensor types change the file ratio. Q4\_0 is also not a claim about
every four-bit checkpoint. Other block types use different
scales, offsets, codebooks or group structures.

The scale in a valid representation can be signed; our decoder
does not force it positive. Selecting a scale and rounding a
floating-point array into codes is a quantizer's job. We constructed
known codes here to make the inverse interpretation unambiguous.
No quantization-quality result follows from reconstructing these
chosen numbers correctly.

\section{From loaded tensors to a ggml computation}
GGUF is a container. ggml is a tensor-computation library. An
inference engine adds a model graph, tokenizer, request state,
generation policy and scheduling around those lower-level pieces.
The ggml authors' introductory tutorial separates tensor objects,
computation graphs, backends and buffers \cite{nguyen2024ggmlintro}.
Keep that separation while tracing the executable example below.

A tensor object records type, dimensions, byte strides and the
location of its data. Creating such an object is not the same
operation as allocating device memory or evaluating a layer.
Our optional reference program loads the trusted fixture through
\texttt{gguf\_init\_from\_file}, retaining both the GGUF directory
context and the ggml context that owns the loaded tensor objects
and data. It then creates a separate work context for inputs,
intermediate objects and computation graphs.

\EngineFigure{foundation-loader-sequence}{UML sequence view of the
optional CPU reference program. Loading establishes tensor objects
and host data; graph construction establishes dependencies; execution
produces results. The ownership lifetimes extend through their last
consumer. This example does not perform GPU transfers or run an LLM.}{fig:foundation-loader-sequence}

Consider the dense tensor. Its first eight floats are
$[1,2,3,4,-1,0,1,2]$. In conventional row-major mathematical
notation, read them as
\[
\mathbf W=
\begin{bmatrix}1&2&3&4\\-1&0&1&2\end{bmatrix}
\in\mathbb R^{2\times4},\qquad
\mathbf x=\begin{bmatrix}1\\2\\4\\8\end{bmatrix}
\in\mathbb R^4.
\]
The output is
\[
\mathbf W\mathbf x=
\begin{bmatrix}1+4+12+32\\-1+0+4+16\end{bmatrix}
=\begin{bmatrix}49\\19\end{bmatrix}.
\]
The GGUF descriptor's dimensions were $[4,2]$, not $[2,4]$.
ggml stores the fastest dimension first: \texttt{ne[0]} is four
and \texttt{ne[1]} is two. For this contiguous F32 tensor,
\texttt{nb[0]} is four bytes and \texttt{nb[1]} is sixteen.
The address of mathematical element $(r,c)$ is
\[
\operatorname{base}+r\,\texttt{nb[1]}+c\,\texttt{nb[0]}.
\]
For $(1,3)$, the file address is $192+16+12=220$.
The four bytes there are \texttt{00 00 00 40}, the little-endian
F32 representation of two. The dimensions change their notation
between the mathematical and library views; the underlying array
has not been transposed.

For a quantized tensor, do not reuse this scalar-address equation
with a fictional half-byte stride. A position identifies a block,
a code within it and the block's metadata. In our Q4\_0 row,
positions zero and sixteen depend on different nibbles of the
same byte. Neither value exists as an independently loadable
four-bit floating-point scalar.

The call \texttt{ggml\_mul\_mat(ctx, weight, input)} creates an
operation node. Under the pinned API's dimension convention,
weights with \texttt{ne=[4,2]} and one input row with
\texttt{ne=[4,1]} produce \texttt{ne=[2,1]}
\cite{ggml2026apipinned}. The program expands a graph from this
output node, then asks the CPU backend to execute that graph
with one thread. Reading the result before execution would be
reading storage, not the result of a completed multiplication.

Here is that boundary in the C++ reference program:
\par\medskip\noindent\begin{minipage}{\linewidth}
\CodeLines{../code/foundations/ggml-reference/main.cpp}{19}{27}
\end{minipage}\par\medskip
The input copy populates the tensor's host data. The shape assertion
checks the operation's contract before execution. Graph expansion
collects the dependencies needed to produce \texttt{y}; computation
then traverses them. Only after the success check does the program
read the result pointer. In this small synchronous CPU path, return
from computation is the completion boundary.

This small path uses host-accessible context-owned data. A device
backend introduces a buffer allocation, a copy or other supported
access mechanism, and an execution-completion boundary before host
readback. A pointer field is not evidence that the CPU may
dereference device storage. The same separation applies when an
engine captures or reuses an execution graph: describing the work
and completing the work remain different events.

The two contexts are kept alive until their last consumers finish.
Destroying the weight context after graph construction but before
execution would invalidate the graph's inputs. Destroying the work
context before reading outputs would invalidate the results.
The example uses C++ ownership wrappers so these lifetimes are
explicit even when a validation check throws an exception.
It reserves a four-MiB work arena for convenience; that number
is not a measured minimum or an engine memory-footprint claim.

\section{Run, compare and deliberately break the byte path}
The default lab needs only Python's standard library:
\begin{verbatim}
python3 code/foundations/gguf_bytes.py
python3 -m unittest discover -s code/foundations -p 'test_*.py'
\end{verbatim}
The first command constructs the fixture in memory and reports its
hash, directory boundary, data origin, tensor intervals and all
decoded values. It writes no file. Predict 256 file bytes, origin
192, and payload intervals $[192,224)$ and $[224,242)$ before
running it. Check those independently before interpreting the values.

To inspect the bytes with an ordinary hex viewer, use
\texttt{--write} followed by a new scratch filename. The command
uses exclusive creation and refuses to overwrite an existing file.
The reverse \texttt{--inspect} command accepts only the documented
teaching subset and reads at most one MiB plus a limit-detection
byte. Do not use it as a general scanner for multi-gigabyte models.
Its deliberately small limits are part of the exercise.

The optional C++ reference route is documented in
\texttt{code/foundations/ggml-reference/README.md}. It builds against
ggml revision \texttt{353b63b}, with the full commit pinned by the
build configuration. It does not download weights. The program
asks ggml to load the fixture independently, verifies its origin,
names, types and dimensions, and executes the dense and packed
projections. The expected dense output is $[49,19]$. The mathematical
packed-row dot product with 32 ones is $-8$, but the actual CPU
path has another rounding boundary, explained next.

\subsection{The first reference run fails for a useful reason}
Our first test demanded that the packed projection equal $-8$ within
$10^{-5}$. It failed: the pinned CPU path returned
$-7.99951171875$. Do not loosen the tolerance before finding the
cause. The direct ggml Q4\_0 decoder reconstructs all 32 weights
exactly. That localizes the discrepancy after weight interpretation,
not in the file offsets or nibble order.

The pinned CPU type table pairs Q4\_0 weights with Q8\_0 activations
for its vector dot product \cite{ggml2026cpupinned}. An input block
of ones receives an ideal Q8\_0 scale $1/127$ and integer codes 127.
The scale is stored as binary16:
\begin{align}
d_x &= \operatorname{round}_{\mathrm{binary16}}(1/127)
       =0.00787353515625,\\
\widehat x &=127d_x=0.99993896484375,\\
\widehat y &=-8\widehat x=-7.99951171875.
\end{align}
This prediction matches the observed value. The file's weight
encoding did not change; the execution path introduced a second
encoding for an operand that arrived as F32. Weight precision,
activation precision and accumulation precision are separate questions.

The finished reference check therefore has two layers. Direct weight
decoding must match the hand values exactly. Graph outputs must match
the prediction including this pinned activation conversion within
$10^{-5}$. Basis-vector inputs receive the same scale correction
at their nonzero coordinate, so the 32 coordinate probes still
catch a permutation. This bound is specific to these fixtures,
not a universal tolerance for quantized models.

The failed initial expectation and the corrected run are recorded in
\texttt{research/measurements/2026-10-03-gguf-byte-lab.md}.
These are observed numerical outputs, not throughput measurements.
Another backend can use a different arithmetic path and require a
different independently justified prediction.

The Python decoder, the hand calculations and ggml's execution path
are distinct checks. They do not share a parser or a dequantization
routine. Agreement therefore gives stronger evidence than writing
and reading a file with the same code. It still establishes only
this fixture and these supported types, not every GGUF model,
every CPU instruction path or every backend.

Now introduce faults one at a time. Move the packed tensor's
relative offset from 32 to zero. The two allocations overlap and
the teaching reader must reject the file. Change its fastest
dimension from 32 to 31: the declared row no longer contains a
complete block. Change the metadata type tag while leaving its
payload unchanged: the field's meaning changes before any tensor
is reached. These failures should name the violated contract,
not appear later as a mysterious matrix-multiplication error.

The test suite tries every truncation of the 256-byte fixture,
including removal of a padding byte. Our canonical teaching
container requires its declared aligned allocation spans to be
present and its padding to be zero. That intentionally strict
policy is stated in the lab; do not quietly apply it to every
third-party artifact without checking the relevant reader's rules.
The tests also cover duplicate names, invalid UTF-8, unsupported
types, enormous dimensions, misalignment and unexplained trailing
bytes.

Finally replace one F32 payload value with a NaN while leaving
all lengths unchanged. Structural inspection succeeds. Numerical
inspection rejects it under the lab's finite-weight policy.
This experiment makes the three contracts visible: correct
container boundaries do not establish acceptable numerical values,
and acceptable values would still not establish a complete
language-model architecture.

\section{What counts as the same model?}
A checksum answers whether bytes match the recorded artifact. It does
not answer whether the runtime interpreted those bytes correctly.
A useful acceptance sequence therefore has more than one level.
First validate container structure. Then validate tensor names, shapes
and encodings against the architecture. Next compare selected numerical
operations. Finally run a known token sequence and compare logits under
the intended tolerance.

Only after the numerical path passes should a text-level test combine
tokenization, templating, model evaluation, sampling and detokenization.
If the combined test fails, retain the smaller tests so you can find
which boundary broke. Looking only at plausible-looking generated text
is a weak test: the wrong model configuration can still produce words.

The same token IDs and weights do not guarantee the same output under
a different sampler or stop policy. Conversely, a sampled string can
change even when the model function is unchanged. Record the numerical
model contract separately from the request's generation policy.
That separation will matter again when caches are shared between requests.

\section{A useful failure is a rejected file}
Run the test suite, which truncates the fixture and gives it an absurd
header length. Both must fail before tensor values are read. Then make
a scratch variant that adds the data origin twice. Predict that element
$(1,2)$ would be read beyond this tiny file rather than at byte 148.
Bounds checks catch this mistake before a language-model test can.

A production loader needs more than this fixture: the complete format
rules, supported versions, duplicate-key handling, resource limits and
tests against maintained implementations. The point of the exercise
is to learn the questions a loader must answer, not to replace a
well-tested format library with a few lines of demonstration code.
You should now be able to draw the byte interval behind a tensor and
explain why loading succeeds or fails before any kernel runs.
\section{Worked problems}
\subsection*{1. Locate a value}
\textbf{Problem.} Where does element $(1,1)$ of the example F32 tensor begin?
\textbf{Solution.} Its row-major index is $1\cdot3+1=4$.
The absolute byte position is $128+4\cdot4=144$; the stored value is 5.

\textbf{Extend the check.} In the GGUF fixture, where is the dense
tensor's mathematical element $(1,3)$? Its data origin is 192, its
descriptor offset is zero and its byte strides are sixteen per row
and four per column. Thus $192+0+1\cdot16+3\cdot4=220$.
The interval $[220,224)$ contains \texttt{00 00 00 40}, which
decodes to F32 value two. Adding the data origin twice would produce
412, outside the file; omitting it would produce 28, inside metadata.
An in-bounds address alone does not establish the correct interpretation.
\subsection*{2. Include the scale}
\textbf{Problem.} Store 64 values in blocks of 32 four-bit codes plus a
16-bit scale per block. How many bytes are needed?
\textbf{Solution.} Each block needs $(32\cdot4+16)/8=18$ bytes.
Two blocks require 36 bytes, not 32.

\textbf{Extend the check.} In the actual Q4\_0 block, which byte
encodes position seventeen? It belongs to the high-nibble half, at
packed-byte index $17-16=1$. The scale begins at absolute byte 224;
two scale bytes precede the packed bytes, so the required byte is
$224+2+1=227$. It contains \texttt{E1}. Its high nibble is fourteen,
giving $0.5(14-8)=3$. The low nibble encodes position one instead,
giving $0.5(1-8)=-3.5$. Reading adjacent nibbles as adjacent
logical coordinates exchanges these roles.

\textbf{Predict the kernel.} The direct decoded-weight sum is $-8$.
For the pinned CPU graph, each input one is reconstructed as
$127\operatorname{round}_{16}(1/127)=0.99993896484375$.
The result is $-7.99951171875$, a difference of $0.00048828125$
from ideal arithmetic. A $10^{-5}$ tolerance against the ideal sum
must fail; the same tolerance against the conversion-aware prediction
passes. Explain the conversion before choosing the comparison.
\subsection*{3. Reject before interpreting}
\textbf{Problem.} Design a truncation test for the teaching fixture.
\textbf{Solution sketch.} Remove its last byte and require the inspector
to reject the payload interval. It must not return a shorter tensor or
silently fill the missing byte with zero.

\textbf{Distinguish the layers.} Truncating the GGUF fixture from
256 to 255 bytes removes allocation padding, not a weight. The lab
still rejects it under its explicit canonical-padding policy. Replacing
the F32 bytes at 220 with a quiet NaN changes no interval: structural
inspection accepts the container, then numerical inspection rejects
the value. Finally, supplying two finite tensors named
\texttt{dense.weight} and \texttt{packed.weight} to a decoder model
does not create its attention, normalization or output parameters.
A model loader must reject missing architecture requirements even
after both earlier checks pass. State which contract each experiment
violates before writing its expected error.

